% =============================================================================
% ANEXO Z — Correos Transaccionales  |  ANEXO AA — Decisiones de Arquitectura
% =============================================================================
\chapter{Anexo Z — Correos Transaccionales}
\label{cap:anexo_correos}

Catálogo de los correos enviados por SMTP (\comando{spring-boot-starter-mail}), asociados a las
operaciones de cuenta y seguridad.\par

\renewcommand{\arraystretch}{1.3}
\fontsize{9}{10.8}\selectfont\sffamily
\begin{tabularx}{\textwidth}{|l|X|l|}
\hline
\rowcolor{colorPrimario}
\textcolor{white}{\textbf{Correo}} & \textcolor{white}{\textbf{Disparador}} & \textcolor{white}{\textbf{Contenido}} \\ \hline
Verificación OTP & Operación sensible & Código de 6 dígitos (15 min). \\ \hline
\rowcolor{grisTecnico}
Recuperación de contraseña & \comando{/forgot-password} & Enlace seguro de restablecimiento. \\ \hline
Confirmación de email & Cambio de email & Enlace de confirmación (\comando{token\_hash}). \\ \hline
\rowcolor{grisTecnico}
Confirmación de baja & Eliminación de cuenta & Aviso de baja con OTP. \\ \hline
Correo de administración & Panel admin & Comunicaciones masivas a usuarios. \\ \hline
\end{tabularx}\normalfont\normalsize

\begin{NotaTecnica}
Para el correo masivo de administración, Tomcat configura \comando{max-http-form-post-size: -1} y
\comando{max-swallow-size: -1} (sin límite), evitando el truncado de formularios extensos.
\end{NotaTecnica}

\clearpage

% =============================================================================
\chapter{Anexo AA — Decisiones de Arquitectura (ADR)}
\label{cap:anexo_adr}

Registro de las decisiones de arquitectura más relevantes del proyecto, con su contexto y
consecuencia, en formato \textit{Architecture Decision Record}.\par

\section{ADR-01: Despliegue Monolítico}
\label{sec:adr_monolito}
\textbf{Contexto:} se requería un despliegue simple y un único artefacto.\par
\textbf{Decisión:} empaquetar la SPA dentro del JAR de Spring Boot (FE en \comando{static/}).\par
\textbf{Consecuencia:} un solo artefacto sirve API y UI; el \comando{SpaFallbackFilter} resuelve
las rutas del cliente. Simplifica el despliegue a costa de acoplar los ciclos de \textit{release}.

\section{ADR-02: Autenticación Delegada en Supabase}
\label{sec:adr_auth}
\textbf{Contexto:} evitar gestionar credenciales y sesiones propias.\par
\textbf{Decisión:} usar Supabase Auth como proveedor de identidad y el backend como
\textit{Resource Server} que valida JWT vía JWKS.\par
\textbf{Consecuencia:} menor superficie de ataque; dependencia de la disponibilidad de Supabase.

\section{ADR-03: jOOQ frente a un ORM}
\label{sec:adr_jooq}
\textbf{Contexto:} se necesitaba control fino del SQL contra el schema \comando{core}.\par
\textbf{Decisión:} adoptar jOOQ (SQL tipado) en lugar de un ORM de propósito general.\par
\textbf{Consecuencia:} consultas explícitas y optimizables; mayor verbosidad que un ORM.

\section{ADR-04: Seguridad por RLS + RPCs SECURITY DEFINER}
\label{sec:adr_rls}
\textbf{Contexto:} aislar los datos por usuario sin lógica duplicada.\par
\textbf{Decisión:} habilitar RLS en \comando{core} y canalizar las escrituras sensibles por RPCs
\comando{SECURITY DEFINER} que usan \comando{auth.uid()}.\par
\textbf{Consecuencia:} aislamiento robusto; las RPCs deben fijar \comando{search\_path}.

\section{ADR-05: Estado del Frontend con Zustand}
\label{sec:adr_zustand}
\textbf{Contexto:} gestionar estado global sin \textit{boilerplate}.\par
\textbf{Decisión:} \textit{stores} Zustand por dominio, con sesión en \comando{sessionStorage}.\par
\textbf{Consecuencia:} sesiones múltiples por pestaña; limpieza vía \comando{resetAllStores}.

\clearpage
